\documentclass[12pt]{article}

\usepackage[francais]{babel}
\usepackage[utf8]{inputenc}
\usepackage[T1]{fontenc}
\usepackage[left=2cm,right=2cm,top=2cm,bottom=2cm]{geometry}
\usepackage{graphicx}

\begin{document}
\begin{titlepage}

\newcommand{\HRule}{\rule{\linewidth}{0.5mm}}

\center

\textsc{\LARGE Université de Rouen}\\[1.5cm]
\textsc{\Large Miniquiz}\\[0.5cm]
\textsc{\large Langage Web 1}\\[0.5cm]

\HRule \\[0.4cm]
{ \huge \bfseries Rapport De Projet}\\[0.2cm] % Titre du document.
\HRule \\[1.5cm]

  \begin{minipage}{0.4\textwidth}
    \begin{flushleft} \large
      \emph{Membres:}\\
      Valentin \textsc{Crochemore}\\
      Yoann \textsc{Fleury}
    \end{flushleft}
  \end{minipage}
  ~
  \begin{minipage}{0.4\textwidth}
    \begin{flushright} \large
      \emph{Enseignant:} \\
      Florent \textsc{Nicart}
    \end{flushright}
  \end{minipage}\\[3cm]


  {\large Date : \today}\\[3cm]

  \includegraphics{res/logo.png}\\[1cm]

  \vfill
\end{titlepage}

\newpage
\tableofcontents
\newpage

\section{Architecture Logicielle}
    \subsection{Base De Données}
        \paragraph{}
        L'architecture de notre base de données est la plus souple possible afin de permettre de futures extensions de l'application. On peut par exemple expliquer notre choix pour les tables suivantes :
        \begin{description}
            \item[mq\_access] Si l'application vient à évoluer, les droits d'accès ne sont pas codés en dur dans le code. Donc pour le moment, cette table contient deux lignes qui définissent l'administrateur et l'utilisateur normal, mais il est possible de rajouter par la suite un nouveau type d'utilisateur et d'ajouter des fonctions pour ceux-ci.
            \item[mq\_answer] Nous avons choisi de stocker les réponses aux questions dans une table différente afin de pouvoir par la suite faire évoluer notre application en proposant plus de 4 réponses à chaque questions. Idem, on pourrait se servir de cette table pour proposer des réponses grâce à des requêtes AJAX et en fonction de la question, des réponses plausibles apparaissent aux yeux du créateur du quiz.
        \end{description}
    
    \subsection{Architecture MVC}

\section{Technologies}
    \subsection{PHP}
        \paragraph{} Nous avons choisi \texttt{PHP} car c'est un langage que nous connaissons bien, libre et qui offre une certaine souplesse. De plus, son grand catalogue de framework nous a motivé à utiliser ce langage. 
    
        \subsubsection{Silex}
            \paragraph{} Silex est un micro-framework \texttt{PHP} basé sur Symfony2 et est disponible sous licence libre MIT. Il est connu pour être concis, extensible et débuggable facilement, ce qui n'est pas forcément le cas de son grand frère. Son API se veut simple d'utilisation et même "fun" à utiliser selon ses créateurs Fabien \textsc{Potencier}, le créateur du framework Symfony et Igor \textsc{Wiedler}.
        
        
        \subsubsection{Composer}
            \paragraph{}Il est rare de nos jours de trouver un projet \texttt{PHP} qui ne contienne pas le gestionnaire de dépendance qu'est Composer. Cet outil libre sous licence MIT permet de gérer des dépendances en PHP et de faciliter l'utilisation des \texttt{namespace} d'un projet grâce à un autoloader généré automatiquement. Cet outil nous a permis d'installer plusieurs dépendances :

            \begin{description}
                \item[Silex] un micro-framework basé sur Symfony
                \item[Doctrine] un ORM qui permet d'ajouter une couche d'abstraction de la base de données pour PHP
                \item[Parsedown] une bibliothèque qui permet d'implémenter une interprétation de la syntaxe Markdown
                \item[Password-Compat] une autre bibliothèque permettant d'assurer la compatibilité des fonctions \texttt{password-hash} et \texttt{passsword-verify} sous une version supérieure ou égale à la version 5.3 de \texttt{PHP} alors que ces dernières fonctions ne sont utilisables que sous une version minimum 5.5 de \texttt{PHP} sans bibliothèque.
                \item[Symfony] un framework très connu dans le monde de \texttt{PHP}. Dans notre cas, nous n'avons pris qu'une version minimale du framework, limité aux outils de débuggage, de routing, de session et de pont entre \texttt{PHP} et \texttt{Twig}.
                \item[Twig] un moteur de templates pour \texttt{PHP}.
            \end{description}
        
            \paragraph{} Toutes les dépendances ci-dessus sont évidemment disponibles sous licence libre et sont réputées pour leur sécurité. Elles sont également toutes disponibles en accès libre sur GitHub pour qui souhaiterait les consulter.

    \subsection{HTML}
        \paragraph{HTML5} HTML5 est la nouvelle norme pour qui souhaite suivre les standards du W3C. De plus, cette version est maintenant bien supporté par les navigateurs, voilà pourquoi nous avons choisi de suivre cette celle-ci. De plus, cette version permet une sémantique plus forte grâce aux nouveaux éléments comme \texttt{header}, \texttt{footer} et \texttt{section}.
        
    \subsection{CSS}
        \paragraph{CSS3} CSS3 est la dernière version de CSS, et nous avons suivi cette nouvelle version dans nos léger développement de feuille de style. Les autres feuilles de style ont été générées pour Bootstrap grâce à SASS.
        
        \paragraph{Bootstrap} Bootstrap est un framework HTML, CSS et Javascript qui permet de rapidement mettre en place un web design dit responsive. C'est à dire que celui-ci va s'adapter en fonction de la taille et de la résolution de l'écran sur lequel il s'affiche. Nous avons choisi ce framework car il permet une grande personnalisation des composants. De plus, sa grande communauté permet d'échanger des astuces pour faire de beaux et bons design rapidement. Il fait parti des grands noms lorsque l'on parle de framework CSS aux côtés de Foundation (qui permet moins de personnalisation simple) et Materialize (qui est toujours en alpha lors de l'écriture de ce rapport).
        
    \subsection{Javascript}
        \paragraph{JQuery} Nous avons choisi d'utiliser JQuery pour la manipulation du DOM. JQuery fait maintenant parti de ces standards dans le développement web, et est très utilisé lorsque l'on doit manipuler l'arbre DOM. Nous l'avons principalement utilisé pour les vérification de mot de passe par exemple.
    
    \subsection{MySQL}
        \paragraph{MySQL} Pour la base de données, quoi de mieux que MySQL. MySQL fait parti de ces conventions de développement web dès que l'on parle de petite application web qui ne demande pas beaucoup de traitement et de ressources.
    

\section{Documentation}

\section{Installation}

\section{Bonne pratique et Sécurité}

\section{Ce que nous avons appris}
    \subsection{Problèmes rencontrés}
    
    
\section{Sigles et terminologie}
    \begin{description}
        \item[CSS] Cascading StyleSheet
        \item[DOM] Document Object Model. Interface indépendante de tout langage de programmation qui permet de manipuler du contenu HTML ou XML.
        \item[HTML] HyperText Markup Language. Langage de balises permettant une représentation sémantique de données.
    \end{description}

\end{document}
